\section{Evaluation Methodology}
\label{sec:evaluation}

The evaluation is a controlled comparison of the native release candidate and
the pinned stock SGLang-Omni comparator on one DGX Spark. It is designed to
answer serving-system questions about latency, throughput, streaming cadence,
resource use, and reliability. It does not evaluate perceptual speech quality.

\subsection{Subjects and workload}

The production matrix contains 12 runs: two implementations, concurrency
profiles B1, B3, and B6, and two rounds. Each run performs 24 unmeasured warm-up
requests followed by at least 200 successfully completed measured requests.
Only one subject may use CUDA during a run. Round one executes stock SGLang
before Native; round two reverses that order and executes Native before stock
SGLang. The Native blocks on either side of the round boundary remain
consecutive, avoiding one unnecessary model reload while preserving the
required round-level order reversal.

Both subjects receive the exact same ordered JSONL workload, model repository
and revision, language, text, VoiceDesign description, seed, portable sampling
controls, streaming mode, and 20.48-s duration ceiling. Workload identity is
bound by SHA-256, including the ordered row seeds. The external Rust HTTP client
and localhost transport are shared.

Native completion is accepted only with explicit natural end-of-sequence,
\code{finish\_reason=stop}, and no length limit. The stock raw-PCM endpoint does
not expose a finish reason; its natural-EOS and length-limit fields must remain
unknown. As a conservative boundary gate, an accepted stock response must be
strictly shorter than 489,600 samples and 20.40 s. Passing that gate excludes a
known length-boundary case but does not prove natural EOS.

\subsection{Client-visible metrics}

The client records a monotonic start time immediately before writing each
request. Time to first audio (TTFA) ends when the first non-empty decoded PCM
payload byte reaches the client. Headers, multipart JSON, WAV headers, framing,
and zero-length transport chunks do not count as audio.

For request $i$, with wall time $w_i$ and decoded audio duration $a_i$,

\begin{equation}
  \mathrm{RTF}_i = \frac{w_i}{a_i}.
\end{equation}

For a concurrent scenario with elapsed wall time $W$ and successful requests
$S$, the throughput-oriented aggregate real-time factor is

\begin{equation}
  \mathrm{RTF}_{\mathrm{aggregate}} =
  \frac{W}{\sum_{i \in S} a_i}.
\end{equation}

The report keeps this quantity distinct from
$\sum_i w_i / \sum_i a_i$, which is a summed-request-wall latency view. It also
reports successful and attempted request throughput, first-to-final PCM
streaming window, packet arrivals, inter-packet gaps, response validity,
timeouts, cancellations, malformed responses, and explicit termination fields.

\subsection{Resource and energy measurement}

Timestamped host, cgroup, NVIDIA, and power telemetry is sampled every 100 ms;
a run fails if any required stream has an observed gap greater than 200 ms.
Measured-phase boundaries exclude warm-up work from performance and energy
reduction. Process RSS and GPU-visible unified memory are reported separately
and are not summed as independent physical memory. Energy is integrated from
the timestamped board-power series after subtraction of a separately measured
idle baseline, then normalized per generated audio minute. An unavailable
sensor remains unavailable rather than becoming an estimate.

\subsection{Fail-closed publication gate}

A run is excluded if another CUDA process is present, evidence is missing or
digest-mismatched, the workload or normalized sampling differs between
subjects, the success minimum is not met, Native completion is not natural EOS,
Stock EOS is imputed, telemetry cadence fails, or the comparator contains an
undisclosed execution change. Setup failures remain in the audit trail but do
not enter distributions. Every accepted evidence path is relative to the
manifest root and bound by byte length and SHA-256.
